% 在这里填写个人信息；英文授予学位名称必须按实际学位填写。
\tongjisetup{
  ctitle={学位论文中文题目},
  etitle={Thesis Title in English},
  cauthor={\blindreviewfield{作者姓名}},
  eauthor={\blindreviewfield{Author Name}},
  studentnumber={\blindreviewfield{学号}},
  cdepartment={\blindreviewfield{学院全称}},
  edepartment={\blindreviewfield{School Name}},
  csubjectcategory={工学},
  esubjectcategory={Engineering},
  cmajorfirst={学科或专业学位类别},
  emajorfirst={Discipline or Degree Category},
  cmajorsecond={二级学科或专业领域},
  emajorsecond={Discipline or Professional Field},
  cresearchfield={研究方向},
  eresearchfield={Research Field},
  csupervisor={\blindreviewfield{导师姓名及职称}},
  esupervisor={\blindreviewfield{Supervisor Name and Title}},
  cassosupervisor={},
  eassosupervisor={},
  cjointinstitution={},
  ejointinstitution={},
  cfunds={},
  efunds={}
  % cdate={二〇二六年十一月}, edate={November 2026},
  % edegree={Master of Engineering}
}
\makeatletter
\tongjideletelinebreakcmd{tongji@ctitle}{tongji@posttitle}
\def\tongji@cheadingtitlenew{\tongji@posttitle}
\makeatother

\begin{cabstract}
本文档是论文排版示例，供填写真实研究内容前检查格式使用。本示例以同济大学研究生院二〇二五年十月发布的写作参考指南为依据，对页面设置、标题层次、摘要、图表、参考文献与附录编号进行了调整。撰写时请用研究目的、研究方法、研究成果和结论替换本段，突出自己的创新成果。摘要内不应出现图表，中文和英文内容应相互对应。
\end{cabstract}
\ckeywords{学位论文, 排版, 格式核对}

\begin{eabstract}
This document is a typesetting example for checking a thesis layout before inserting research content. It follows the October 2025 writing guide issued by the Graduate School of Tongji University. The layout covers page dimensions, headings, abstracts, figures, tables, references, and appendix numbering. Replace this paragraph with the purpose, methods, results, and conclusions of the research. The English abstract should correspond to the Chinese abstract.
\end{eabstract}
\ekeywords{thesis, typesetting, format verification}
